% =====================================================================
%  Market Regime Detection Using Autoencoders and PCA
%  Presentation Script (15 min) + Q&A Preparation
%  Compile with pdflatex (run twice for the table of contents).
%  Only standard packages are used (works on Overleaf out of the box).
% =====================================================================
\documentclass[11pt]{article}

\usepackage[a4paper,margin=2.2cm]{geometry}
\usepackage[T1]{fontenc}
\usepackage[utf8]{inputenc}
\usepackage{lmodern}
\usepackage{microtype}
\usepackage{xcolor}
\usepackage{titlesec}
\usepackage{enumitem}
\usepackage{fancyhdr}
\usepackage{array}
\usepackage{booktabs}
\usepackage[most,skins,breakable]{tcolorbox}
\usepackage{hyperref}

% ---------- Colour palette -------------------------------------------
\definecolor{ink}{HTML}{1A1A2E}
\definecolor{soft}{HTML}{F4F6FB}
\definecolor{rule}{HTML}{D7DEEA}

% Person 1 — Data (blue)
\definecolor{p1}{HTML}{1565C0}
\definecolor{p1l}{HTML}{E3F0FC}
% Person 2 — Autoencoder (teal/green)
\definecolor{p2}{HTML}{00897B}
\definecolor{p2l}{HTML}{DEF3F0}
% Person 3 — PCA (purple)
\definecolor{p3}{HTML}{6A1B9A}
\definecolor{p3l}{HTML}{F1E6F8}
% Accents
\definecolor{amber}{HTML}{E08A00}
\definecolor{amberl}{HTML}{FBEFD6}
\definecolor{crimson}{HTML}{C62828}

% current accent (switched per section)
\colorlet{accent}{p1}
\colorlet{accentl}{p1l}

% ---------- Section styling ------------------------------------------
\titleformat{\section}
  {\sffamily\LARGE\bfseries\color{accent}}{}{0pt}
  {}[{\color{accent}\titlerule[2pt]}]
\titlespacing*{\section}{0pt}{18pt}{8pt}

\titleformat{\subsection}
  {\sffamily\large\bfseries\color{accent}}{}{0pt}{}
\titlespacing*{\subsection}{0pt}{12pt}{4pt}

\setlength{\parindent}{0pt}
\setlength{\parskip}{4pt}
\renewcommand{\familydefault}{\sfdefault}

% ---------- Header / footer ------------------------------------------
\pagestyle{fancy}
\fancyhf{}
\renewcommand{\headrulewidth}{0.6pt}
\renewcommand{\headrule}{\hbox to\headwidth{\color{rule}\leaders\hrule height \headrulewidth\hfill}}
\fancyhead[L]{\footnotesize\color{ink}\textbf{Market Regime Detection} — AE \& PCA}
\fancyhead[R]{\footnotesize\color{accent}\leftmark}
\fancyfoot[C]{\footnotesize\color{ink}\thepage}

\hypersetup{colorlinks=true, linkcolor=accent, urlcolor=p1}

% ---------- Reusable boxes -------------------------------------------
% Timing cue box: \begin{cue}{label}{duration} ... \end{cue}
\newtcolorbox{cue}[2]{
  enhanced, breakable, sharp corners=downhill,
  colback=accentl, colframe=accent, boxrule=0pt,
  leftrule=4pt, top=5pt, bottom=5pt, left=9pt, right=9pt,
  before skip=8pt, after skip=8pt,
  overlay unbroken and first={
    \node[anchor=north east, font=\sffamily\bfseries\footnotesize,
          text=white, fill=accent, inner sep=3pt, rounded corners=2pt]
      at ([xshift=-6pt,yshift=-5pt]frame.north east) {#2};},
  title={\sffamily\bfseries\color{accent}#1},
  fonttitle=\sffamily\bfseries
}

% Q&A box: \begin{qa}{question} answer \end{qa}
\newtcolorbox{qa}[1]{
  enhanced, breakable, sharp corners,
  colback=white, colframe=accent, boxrule=1pt,
  left=10pt, right=10pt, top=8pt, bottom=8pt,
  before skip=7pt, after skip=7pt,
  coltitle=white, colbacktitle=accent,
  title={\sffamily\bfseries Q.\hspace{4pt}#1}
}

% Key-numbers card
\newtcolorbox{numbers}[1]{
  enhanced, sharp corners,
  colback=soft, colframe=accent, boxrule=0pt, leftrule=4pt,
  left=10pt, right=10pt, top=6pt, bottom=6pt,
  title={\sffamily\bfseries\color{accent}#1}, fonttitle=\sffamily\bfseries
}

% Person banner
\newcommand{\personbanner}[3]{%
  \begin{tcolorbox}[enhanced, sharp corners, colback=accent, colframe=accent,
      left=14pt, right=14pt, top=10pt, bottom=10pt, before skip=4pt, after skip=10pt]
    {\sffamily\bfseries\LARGE\color{white} #1}\\[2pt]
    {\sffamily\color{white}\large #2}\hfill{\sffamily\bfseries\color{white} #3}
  \end{tcolorbox}}

\newcommand{\setp}[2]{\colorlet{accent}{#1}\colorlet{accentl}{#2}}

% =====================================================================
\begin{document}

% ---------------- Title page -----------------------------------------
\begin{titlepage}
\setp{p1}{p1l}
\centering
\vspace*{2cm}
{\sffamily\bfseries\Huge\color{ink} Market Regime Detection}\\[6pt]
{\sffamily\bfseries\Huge\color{p1} Using Autoencoders \& PCA}\\[18pt]
{\color{rule}\rule{0.7\textwidth}{1.5pt}}\\[18pt]
{\sffamily\Large\color{ink} Presentation Script \;\&\; Q\&A Preparation}\\[6pt]
{\sffamily\large\color{ink} 15-minute group talk \,$\cdot$\, three speakers \,$\cdot$\, 5 minutes each}\\[40pt]

\begin{tcolorbox}[enhanced, width=0.8\textwidth, halign=center,
   colback=soft, colframe=rule, boxrule=1pt, sharp corners]
\sffamily\large
\textcolor{p1}{\textbf{Person 1 — Harjot Singh}}\\ Data Collection \& Preprocessing\\[8pt]
\textcolor{p2}{\textbf{Person 2 — Eleni Tsaousi}}\\ Autoencoder \& Clustering\\[8pt]
\textcolor{p3}{\textbf{Person 3 — Adene Dinoshi}}\\ PCA, Comparison \& Interpretation
\end{tcolorbox}

\vfill
{\sffamily\small\color{ink} Investments — Selected Quantitative Tools, University of Zurich \\ Supervisor: Matthias Feiler}
\end{titlepage}

% ---------------- Table of contents ----------------------------------
\setp{ink}{soft}
{\hypersetup{linkcolor=ink}
\renewcommand{\contentsname}{\sffamily\color{ink}Contents}
\tableofcontents}
\vspace{6pt}

\begin{tcolorbox}[enhanced, sharp corners, colback=amberl, colframe=amber,
    boxrule=0pt, leftrule=4pt, left=10pt, right=10pt, top=6pt, bottom=6pt]
\sffamily\small
\textbf{\textcolor{amber}{How to use this document.}} Part~I is the word-for-word
script, split into three colour-coded sections — one per speaker — with a timing
cue and a ``key numbers'' card on every block. Part~II is the Q\&A: each person
should master their own colour-coded section and skim the others. All figures
referenced live in \texttt{pca-regime-analysis/outputs/figures/}.
\end{tcolorbox}

\newpage

% #####################################################################
% PART I — SCRIPT
% #####################################################################
\setp{ink}{soft}
\begin{center}
{\sffamily\bfseries\Huge\color{ink} Part I \;\textemdash\; Presentation Script}
\end{center}
\vspace{4pt}

% =====================================================================
% PERSON 1
% =====================================================================
\setp{p1}{p1l}
\section{Person 1 — Data Layer}
\personbanner{Harjot Singh}{Data Collection \& Preprocessing}{$\approx$ 5 min}

\begin{numbers}{Key numbers to remember}
\sffamily\small
Weekly data, \textbf{Apr 1988 → May 2026} ($\approx$38 yrs) \;$\cdot$\;
\textbf{30 features} (8 log-return + 22 macro) \;$\cdot$\;
Split \textbf{70/15/15 by time} \;$\cdot$\;
Window \textbf{26 weeks}, stride 1 \;$\cdot$\;
Train windows \textbf{1365 $\times$ 26 $\times$ 30} \;$\cdot$\;
\textbf{13/13} contract tests pass.
\end{numbers}

\begin{cue}{Opening}{30 sec}
Good morning, everyone. My name is Harjot, and I will start us off by walking you
through the foundation of our project: where our data comes from, why we chose it,
and how we turned raw market information into a clean, model-ready dataset that
both my teammates could build on.
\end{cue}

\begin{cue}{Why this data?}{60 sec}
The central question we study is: can we detect distinct \textbf{market regimes} —
recurring states like bull markets, crisis periods, or high-inflation environments —
from financial and macroeconomic data alone, with \emph{no} manual labelling?

To do that we needed data capturing both price dynamics \emph{and} the economic
environment. We chose a weekly US dataset from April 1988 to May 2026 — nearly
38 years. It has 30 features: eight price/index series (equity market, gold, oil,
the US dollar, Treasury yields) plus 22 macro and valuation indicators — yield
curve, CPI, unemployment, CAPE, dividend yield, GDP, recession probability.
Together these tell us \emph{what} markets are doing and \emph{why} — exactly what
regime detection needs.
\end{cue}

\begin{cue}{The stationarity problem}{60 sec}
My first problem was non-stationarity. The eight price columns are raw levels —
they just trend upward over decades. Fed in directly, the autoencoder would learn
the \emph{trend}, not the \emph{regime}. Augmented Dickey–Fuller tests confirmed
all eight had p-values near 1.0 — strongly non-stationary. The fix: convert each
to \textbf{log returns} — the natural log of today's price over yesterday's. After
that, all eight had p-values below 0.0001. The 22 macro features were already
stationary — year-over-year GDP or unemployment don't trend forever — so I kept
them as-is.
\end{cue}

\begin{cue}{Split \& scaling — anti-leakage}{60 sec}
The second critical step was preventing \textbf{data leakage}. We split 70/15/15
\emph{by time} — never shuffled. Train is 1988–2014, validation 2014–2020, and the
test set is the most recent period, 2020–2026. I fit a \texttt{StandardScaler} on
the \emph{training data only}, then applied it to validation and test. The model
must never see the future when learning to standardise.

I also wrote a 13-point automated contract test: no NaNs, correct shapes and
dtypes, and — crucially — that validation and test means \emph{drift away} from
zero, proving the scaler was not re-fit on them. All 13 checks pass.
\end{cue}

\begin{cue}{Windows \& the handoff}{60 sec}
Finally, the input windows. A sequence autoencoder needs a fixed-length snapshot
of the recent past. I tested windows from 13 to 78 weeks and measured the regime
\emph{flip rate} between adjacent windows — it stayed under 0.5\% for every size,
because slow macro signals dominate. So I chose \textbf{26 weeks} ($\approx$ 6
months) on economic grounds — a natural half-year horizon. The training tensor is
$1365 \times 26 \times 30$: 1365 overlapping windows, each 26 weeks of 30 features.

One honest caveat we flag in the report: because the windows overlap by 25 of 26
weeks, those 1{,}365 windows are far fewer \emph{independent} economic episodes
than the count suggests — Adene will return to what that means for our results.
The output is a clean, leak-free, reproducible dataset that both my teammates
built directly on.
\end{cue}

\begin{cue}{Hand over}{10 sec}
I'll now hand over to Eleni, who took this data and trained the autoencoder.
\end{cue}

% =====================================================================
% PERSON 2
% =====================================================================
\setp{p2}{p2l}
\section{Person 2 — Autoencoder}
\personbanner{Eleni Tsaousi}{Autoencoder Implementation, Training \& Clustering}{$\approx$ 5 min}

\begin{numbers}{Key numbers to remember}
\sffamily\small
Input \textbf{780} ($26\times30$) \;$\cdot$\;
Encoder \textbf{780→256→64→16} (ReLU) \;$\cdot$\;
Latent \textbf{16-D} \;$\cdot$\;
Adam, weight decay \textbf{5e-4}, early stopping (patience 40), $\le$200 epochs \;$\cdot$\;
Recon loss \textbf{train 0.30 / val 0.61 / test 1.20} \;$\cdot$\;
KMeans \textbf{K=4} (fit on train latent only).\\[3pt]
Ablation: 16-D gives \textbf{lowest test loss}; OOS regime collapse persists at every width.\\[3pt]
Regime counts (full data): \textbf{R0 = 379, R1 = 666, R2 = 345, R3 = 521}.
\end{numbers}

\begin{cue}{Opening}{20 sec}
Thank you, Harjot. I'm Eleni, and I built and trained the autoencoder that learns
a compressed representation of market behaviour from the windows Harjot produced.
\end{cue}

\begin{cue}{Why an autoencoder?}{50 sec}
An autoencoder is a neural network that compresses its input through a narrow
\textbf{bottleneck} and then reconstructs it. The bottleneck forces it to keep only
what's essential — the underlying market state, not the noise. Its key advantage
over PCA, which Adene will cover, is capturing \textbf{non-linear} relationships.
Markets aren't linear: a 1\% rise in yields means something very different
depending on whether inflation is high or low and whether unemployment is rising.
We want a representation that can pick up those interactions.
\end{cue}

\begin{cue}{Architecture}{60 sec}
The input is one flattened 26-week window: $26 \times 30 = 780$ values. The encoder
compresses through three layers — 256, then 64, then a \textbf{16-dimensional
bottleneck} — with ReLU activations. The decoder mirrors that back up to 780. We
reconstruct the window and use \textbf{mean squared error} as the loss.

Why 16 dimensions? The data analysis showed a 10-component PCA already explains
67\% of variance and many features are strongly correlated — 11 pairs above 0.8 —
so the 30 inputs carry far fewer independent dimensions. We didn't just assume 16,
though: we justified it with an ablation, which I'll show in a moment.
\end{cue}

\begin{cue}{Training, early stopping \& reconstruction}{55 sec}
I trained with the \textbf{Adam} optimiser and weight decay of $5\times10^{-4}$,
and used \textbf{early stopping} on the validation loss — up to 200 epochs, keeping
the best checkpoint. This is our main defence against overfitting the training
period.

Reconstruction loss reached about \textbf{0.30} on train, \textbf{0.61} on
validation, and \textbf{1.20} on test. That jump on test is the first important
result of the paper: it reflects genuine \textbf{distributional shift}. The test
period, 2020–2026, includes COVID, the post-pandemic inflation surge, and the
sharpest Fed hiking cycle in decades — conditions barely present in training. We
read it as evidence of a regime shift, not a model bug.
\end{cue}

\begin{cue}{Latent-dimension ablation — why 16, and an honest finding}{60 sec}
To justify the bottleneck width we retrained the model at 2, 3, 8, 16 and 32
dimensions. Two findings, both in the paper. First, reconstruction improves up to
about eight dimensions and then plateaus, and \textbf{16 gives the lowest
out-of-sample reconstruction loss} — going wider to 32 actually overfits slightly.
That's our reason for choosing 16. Second, and more interesting: the
out-of-sample \emph{regime} coverage does \textbf{not} improve with capacity — the
validation period collapses to only two of the four regimes whether the bottleneck
is 2- or 32-dimensional. So the collapse isn't a capacity problem we can engineer
away; it's the data itself shifting out of sample. That honest result sets up
everything Adene discusses next.
\end{cue}

\begin{cue}{Clustering into regimes}{50 sec}
Once trained, I extracted the 16-D latent vector for every window across all
splits, then applied \textbf{K-Means with K=4}, fitting the cluster centres on the
\emph{training} latent vectors only and assigning regimes to validation and test.
Four clusters matches the intuition that markets move through a few broad states:
expansion, contraction, stress, and transition.

Across the full dataset the regimes split as \textbf{379 / 666 / 345 / 521}
windows. Regime~1 is the largest, common ``baseline'' state; the smaller regimes
pick out the more distinctive, atypical periods.
\end{cue}

\begin{cue}{Hand over}{20 sec}
I exported the latent vectors, regime labels, and reconstruction losses to CSV.
Adene used these as her baseline for comparison. Over to her.
\end{cue}

% =====================================================================
% PERSON 3
% =====================================================================
\setp{p3}{p3l}
\section{Person 3 — PCA \& Interpretation}
\personbanner{Adene Dinoshi}{PCA, Comparison \& Market Regime Interpretation}{$\approx$ 5 min}

\begin{numbers}{Key numbers to remember}
\sffamily\small
\textbf{16-component} PCA (fit on train) \;$\cdot$\;
Variance: 10 PC = 67\%, \textbf{16 PC $\approx$ 71\%}, 50 PC = 82\% \;$\cdot$\;
\textbf{ARI = 0.52}, \textbf{NMI = 0.58} (moderate agreement) \;$\cdot$\;
PCA train regimes: R0 26.6\%, R1 34.0\%, R2 27.3\%, \textbf{R3 12.1\% (crisis)} \;$\cdot$\;
Trading Sharpe: \textbf{train 0.23 / val 0.13 / test $-$0.28}.\\[3pt]
Signal-set ablation (macro / risk / full): \textbf{none positive OOS} — failure is structural, not feature choice.
\end{numbers}

\begin{cue}{Opening}{20 sec}
Thank you, Eleni. I'm Adene. My role was to build a PCA-based regime detector,
compare it rigorously with the autoencoder, produce the visualisations, and
interpret what the regimes mean economically.
\end{cue}

\begin{cue}{Why PCA as a baseline?}{50 sec}
PCA — Principal Component Analysis — is the natural \textbf{linear benchmark}. It
finds the directions of maximum variance and projects the data onto them. It's
fast, interpretable, and well understood. If the autoencoder doesn't clearly beat
PCA, we can't justify the extra complexity. To keep the comparison fair, I used the
\emph{same} 16-dimensional projection as Eleni's bottleneck, fit on training only,
and clustered with the same K-Means, K=4.
\end{cue}

\begin{cue}{What PCA explains}{50 sec}
The scree plot shows 16 components explain about \textbf{71\%} of the variance. For
context: 10 components give 67\%, and you need 50 to reach 82\%. So the 16-D linear
representation already retains most of the information — meaning both methods work
from a genuinely compressed view of the same data. Differences in regime quality
therefore reflect \emph{linear vs.\ non-linear capacity}, not a difference in
dimensionality.
\end{cue}

\begin{cue}{Agreement metrics — the key result}{60 sec}
I measured how much the two methods agree with two standard metrics. The
\textbf{Adjusted Rand Index} is \textbf{0.52} out of 1; \textbf{Normalised Mutual
Information} is \textbf{0.58}. Both mean \emph{moderate} agreement — well above
random, well short of perfect. This is our key finding: PCA and the autoencoder
share real regime structure but disagree on roughly half the boundary placements.
The autoencoder is capturing non-linear structure that linear PCA cannot reproduce
even in 16 dimensions.
\end{cue}

\begin{cue}{Economic interpretation}{60 sec}
The per-regime feature profiles give economic meaning. \textbf{Regime~0}: low
short- and long-term rates — a low-volatility, easing environment. \textbf{Regime~2}:
high CPI and rising 10-year yields — inflationary pressure. \textbf{Regime~3} is the
most distinct — highest yield-spread stress and unemployment with the sharpest GDP
contraction: our \textbf{crisis} regime, about 12\% of windows. \textbf{Regime~1} is
the largest and most diffuse — a normal/transitional baseline. The timeline confirms
this: Regime~3 clusters around the early-1990s recession, the 2008 crisis, and parts
of the COVID shock.
\end{cue}

\begin{cue}{Do the regimes pay? The trading test}{55 sec}
Interpretable regimes are nice, but do they carry \emph{economic} value? We built a
deliberately simple test: using \textbf{training data only}, rank the regimes by
their average one-week-forward market return, then go long the best regime, short
the worst, flat the rest. We freeze that rule and apply it untouched to validation
and test. In sample it works — Sharpe \textbf{0.23} — and it's economically
coherent: the high-stress regime precedes the worst returns. But out of sample it
falls apart: validation Sharpe \textbf{0.13}, and on test \textbf{$-$0.28}, losing
to a flat buy-and-hold benchmark, with an information coefficient indistinguishable
from zero.
\end{cue}

\begin{cue}{Is it the features? Signal-set experiment}{50 sec}
A natural objection is that we simply picked the wrong features. So we ran a
\textbf{signal-set ablation}: we retrained the whole pipeline three times — on a
10-feature \textbf{macro} set, a 10-feature \textbf{risk-and-sentiment} set, and
the full 30 features. The result is striking: \textbf{none} of the three achieves a
positive out-of-sample Sharpe. The risk/sentiment set even reconstructs the test
period best yet trades the worst. This tells us the generalisation failure is a
\textbf{structural property of the post-2020 period}, not an artefact of feature
choice — which is exactly what Eleni's latent ablation also implied.
\end{cue}

\begin{cue}{Limitations \& what we'd improve}{55 sec}
We're upfront about the limits. The windows overlap heavily — 26-week windows one
week apart — so our 1{,}365 windows are far fewer \emph{independent} episodes,
which inflates apparent stability. The autoencoder is trained only to
\emph{reconstruct}; nothing in its objective rewards separating future returns, so
good regimes are a by-product, not a target. And the agreement with PCA is only
moderate — much of the structure is linear, so the neural network's added value is
real but modest. To improve it, the report points to three concrete directions:
\textbf{retrain on a rolling window} so the regimes adapt instead of freezing in
2014; add a \textbf{clustering-aware term} to the loss so the latent space is built
for regime separation; and add \textbf{transaction costs and position sizing}
before treating any of this as tradeable.
\end{cue}

\begin{cue}{Conclusion}{35 sec}
So, what did the study find? An autoencoder \emph{can} learn interpretable market
regimes from unlabelled data, and it captures some non-linear structure a linear
PCA of equal size cannot — but the gain over PCA is modest, and the regimes do not
yet generalise across a genuine structural break. The honest conclusion of the
paper is that this is a useful \textbf{descriptive lens} on market history, not a
ready-to-deploy trading signal — and, importantly, we showed \emph{why}, through
the reconstruction jump, the regime collapse, and two independent ablations. Thank
you — we're happy to take questions.
\end{cue}

\newpage
% #####################################################################
% PART II — Q&A
% #####################################################################
\setp{ink}{soft}
\begin{center}
{\sffamily\bfseries\Huge\color{ink} Part II \;\textemdash\; Q\&A Preparation}
\end{center}
\vspace{2pt}
\begin{tcolorbox}[enhanced, sharp corners, colback=amberl, colframe=amber,
    boxrule=0pt, leftrule=4pt, left=10pt, right=10pt, top=6pt, bottom=6pt]
\sffamily\small \textbf{\textcolor{amber}{Tip.}} Know your own section cold; skim the
other two so you can field a hand-off question gracefully. If you're unsure, it's
fine to say ``that's more Eleni's area'' and redirect.
\end{tcolorbox}

% --------------------- Person 1 Q&A ----------------------------------
\setp{p1}{p1l}
\section{Q\&A — Person 1 (Data)}

\begin{qa}{Why weekly frequency, not monthly or daily?}
Weekly has enough resolution to catch transitions as they develop — monthly might
already be mid-transition — while avoiding daily microstructure noise. Macro series
like CPI and GDP are monthly/quarterly, so we forward-fill them to weekly, which is
fine for slow-moving signals.
\end{qa}

\begin{qa}{Why log returns instead of simple percentage returns?}
Log returns are additive over time, approximately normal for small moves, and the
standard in financial econometrics. The numerical difference is tiny at this scale,
but it makes the ADF stationarity tests cleaner and is methodologically correct.
\end{qa}

\begin{qa}{You split 70/15/15 by time — why not k-fold cross-validation?}
K-fold shuffles time, so the model would train on the future and validate on the
past — forward-looking bias that invalidates any out-of-sample claim. Financial
data is time-ordered; forward chaining (train oldest → test newest) honestly
simulates production.
\end{qa}

\begin{qa}{Why a 26-week window? Would a different size change the regimes?}
We tested 13, 26, 39, 52, 65 and 78 weeks and measured the regime flip rate between
adjacent windows. It stayed near 0.5\% for every size, because slow macro variables
dominate. We picked 26 weeks ($\approx$ 6 months) as a natural economic half-year
rather than over-fitting the window to the data.
\end{qa}

\begin{qa}{What would your 13 contract tests catch if someone introduced leakage?}
The decisive check is that validation/test means \emph{drift away} from zero. If the
scaler had been fit on all data, every split would have mean $\approx 0$,
std $\approx 1$. The test requires the validation mean to \emph{not} be near zero —
proving only training data fit the scaler.
\end{qa}

\begin{qa}{Which 8 columns become log returns, and why those?}
The eight price/index \emph{levels}: equity market (\texttt{\_MKT}), gold
(\texttt{\_AU}), US dollar (\texttt{\_DXY}), corporate bonds (\texttt{\_LCP}),
Treasury (\texttt{\_TY}), oil (\texttt{\_OIL}), value (\texttt{\_VA}), growth
(\texttt{\_GR}). They trend, so they're non-stationary as levels. The other 22 are
already stationary rates/ratios (yields, CPI, CAPE, unemployment\ldots) so we keep
them as-is.
\end{qa}

\begin{qa}{Could you have added more features?}
Yes, but we deliberately kept only series available continuously over the full
1988–2026 span to avoid survivorship bias, and more inputs would enlarge the
780-dimensional vector and make the AE's compression harder without obviously
adding regime-relevant signal.
\end{qa}

% --------------------- Person 2 Q&A ----------------------------------
\setp{p2}{p2l}
\section{Q\&A — Person 2 (Autoencoder)}

\begin{qa}{Why a fully-connected autoencoder, not an LSTM or Transformer?}
The input is a fixed-length 26-week window flattened to 780 values. We want the
overall \emph{state} of that window, not a next-step prediction, so order inside the
window matters less. A dense AE is a natural, efficient baseline; an LSTM or
Transformer adds complexity we haven't yet shown is needed to beat it.
\end{qa}

\begin{qa}{Why 16 latent dimensions?}
Two reasons. First, a 10-component PCA already explains 67\% of variance and 16
explains $\approx$ 71\%, so the effective dimensionality is well below 30. Second,
an earlier 3-D bottleneck showed \emph{mode collapse} — the encoder pushed
everything into one region, giving unstable clusters. Widening to 16 fixed that, and
matching the PCA dimension keeps the comparison fair.
\end{qa}

\begin{qa}{What is early stopping and why does it matter here?}
It monitors validation loss each epoch and keeps the best checkpoint, stopping after
40 epochs with no improvement. Without it the model keeps memorising the training
set — training loss falls while validation loss rises. Since we want regime
structure that generalises across time, we select the model that's best on held-out
data.
\end{qa}

\begin{qa}{Why is the test reconstruction loss (1.20) so much higher than train (0.30)?}
The test period (2020–2026) includes COVID, the fastest inflation in 40 years, and
525 bp of Fed hikes in 18 months — combinations the training era (1988–2014) never
saw together. Higher reconstruction loss means those windows are genuinely novel
relative to training. We read it as a regime shift, not a failure.
\end{qa}

\begin{qa}{Why K-Means with K=4? Did you try other K?}
Honestly, the cluster-validity metrics don't single out K=4 --- silhouette rises
monotonically from about 0.25 at K=2 to 0.35 at K=8, and the inertia elbow is
shallow. So K=4 is chosen on \emph{economic} grounds: four states map cleanly to
easing, inflationary, stress, and a transitional baseline. We deliberately trade a
little silhouette for interpretability and out-of-sample stability, because larger
K fragments into clusters that are harder to interpret and less stable OOS.
\end{qa}

\begin{qa}{Why fit K-Means only on training, then predict val/test?}
Same logic as scaling: we can't use the future to define regimes. Clustering on all
data would let val/test points move the centroids, leaking future information into
the labels. Fitting on training only makes val/test assignments true out-of-sample
predictions.
\end{qa}

\begin{qa}{Regime 1 has 666 windows — far more than the others. Is that a problem?}
No — Regime~1 is the common ``baseline'' state, and markets spend most of their time
in normal conditions, so a large baseline cluster is expected. What matters is that
the smaller, distinctive regimes (e.g.\ the crisis state) are coherent and stable
over time, which they are.
\end{qa}

\begin{qa}{Could this drive a live trading signal today?}
Not yet. The backtest is profitable and sensible in sample but the test Sharpe is
negative ($-0.28$) and IC near zero — consistent with the high test reconstruction
loss. You'd first need to handle the distribution shift, e.g.\ periodic retraining
or online updating.
\end{qa}

% --------------------- Person 3 Q&A ----------------------------------
\setp{p3}{p3l}
\section{Q\&A — Person 3 (PCA \& Interpretation)}

\begin{qa}{What does ``ARI = 0.52'' actually mean?}
The Adjusted Rand Index measures how often both methods agree on whether two windows
are in the same regime, corrected for chance: 1 = perfect, 0 = random. At 0.52 they
share real structure but disagree on about half the cases — where they agree, they
capture the same signal; where they differ, the AE's non-linearity groups things
differently than linear PCA.
\end{qa}

\begin{qa}{Is 71\% variance in 16 components ``enough''?}
For regime detection it appears sufficient — the agreement heatmap shows each PCA
regime has a clear majority overlap with an AE regime, so they share the dominant
structure. The discarded $\approx$ 29\% may hold the non-linear interactions the AE
uses to refine boundaries; whether that refinement is economically meaningful is
exactly what the ARI/NMI quantify.
\end{qa}

\begin{qa}{How do you know the regime labels map to real economic states?}
From the mean standardised feature values per regime. A regime with high CPI and
rising 10-year yields, whose windows cluster around known inflationary periods, is a
consistent story. We cross-checked the crisis regime against recession dates and
found strong overlap. These are reasoned inferences, not ground truth — none exists
for unlabelled regimes.
\end{qa}

\begin{qa}{Why does the test period look dominated by one or two regimes?}
2020–2026 is exceptional — COVID, the reopening boom, the sharpest hiking cycle in
decades. The training era saw those ingredients separately but never in this
sequence, so the model may be pushing genuinely novel windows into the nearest
training-era cluster rather than naming a new state.
\end{qa}

\begin{qa}{Why use t-SNE for the 2-D/3-D plots instead of just PC1 vs.\ PC2?}
Both representations are 16-dimensional. Plotting only the first two or three axes
would throw away most of the information and could exaggerate or hide separation.
t-SNE collapses all 16 dimensions to 2 or 3 while preserving local neighbourhood
structure — a more honest picture of how regimes actually sit in the latent space.
\end{qa}

\begin{qa}{Why interpret regimes from the last week of each window, not the average?}
The last week is the point in time the label is attached to. Averaging all 26 weeks
would smear in early-window events no longer relevant at the window's end. The
last-week snapshot captures the market state at the moment of classification.
\end{qa}

\begin{qa}{If you had to name one key takeaway, what is it?}
The ARI of 0.52. It confirms non-linear representation learning finds \emph{different}
regime structure than the linear baseline — validating the core hypothesis — but the
substantial overlap also shows a large part of the signal is linear and doesn't need
a neural network. Honest verdict: the AE adds something, but not everything, over PCA.
\end{qa}

\begin{qa}{A ``mixed/transitional'' regime sounds vague — is it useful?}
Yes, as a null category. Forcing every window into a sharply named regime risks
over-interpreting noise. A diffuse transitional cluster keeps the other three sharp
and coherent; those windows are real periods of the market moving between states.
\end{qa}

\begin{qa}{Maybe the strategy just used the wrong features — did you check?}
Yes, that's exactly the signal-set experiment. We retrained the entire pipeline on
three feature sets: a 10-feature macro set, a 10-feature risk-and-sentiment set,
and the full 30 features. None achieved a positive out-of-sample Sharpe. Tellingly,
the risk/sentiment set reconstructed the test period best yet traded the worst. So
the failure is not about feature selection --- it's a structural break in the
post-2020 test period that no feature subset escapes.
\end{qa}

\begin{qa}{What single change would most likely make this generalise?}
Rolling retraining. Right now the autoencoder and the K-Means centroids are frozen
at the 2014 training boundary, so a 2020s window is forced onto stale 1990s--2010s
structure. Periodically refitting on an expanding or rolling window would let the
latent space and the regimes adapt as the economy changes. After that, adding a
clustering- or prediction-aware term to the loss so the latent space is optimised
for something other than pure reconstruction.
\end{qa}

% --------------------- Cross-team Q&A --------------------------------
\setp{ink}{soft}
\section{Q\&A — Cross-Team Questions}
\setp{amber}{amberl}

\begin{qa}{How does any of this connect to real investment decisions?}
We built a simple proof-of-concept: assign long/short/flat by which regime
historically preceded the best or worst forward returns. It's profitable in sample
but doesn't generalise — negative test Sharpe ($-0.28$). We're explicit that, as it
stands, regime detection here is a research tool, not a tradeable signal.
\end{qa}

\begin{qa}{Why not Hidden Markov Models, the classic regime method?}
HMMs assume a Markov transition structure and need the number of states and emission
distributions specified up front. Autoencoders learn the representation purely from
data without those assumptions — which is precisely the linear-vs-nonlinear
comparison this project was asked to make. HMMs are a natural extension.
\end{qa}

\begin{qa}{Is 38 years of data enough to trust the results?}
It spans multiple business cycles, two major bear markets (2000–02, 2008–09), two
inflation episodes, and several credit-stress events — reasonable for studying
multi-year macro regimes. It is \emph{not} long enough for strong claims about very
rare tail events, which we don't make.
\end{qa}

\vspace{10pt}
\begin{center}
{\color{rule}\rule{0.6\textwidth}{1pt}}\\[6pt]
{\sffamily\small\color{ink} End of document \;$\cdot$\; good luck with the presentation!}
\end{center}

\end{document}
